\documentclass[12pt]{article}

\usepackage[spanish, es-noshorthands]{babel}
\usepackage[a4paper, margin=3cm]{geometry}
\usepackage[T1]{fontenc}
\usepackage[utf8]{inputenc}
\usepackage{xcolor}
\usepackage{amsmath}
\usepackage{graphicx}
\usepackage{enumitem}
\usepackage{hyperref}
\usepackage{listings}
\usepackage{caption}

\definecolor{codebg}{RGB}{248,248,248}
\definecolor{codeframe}{RGB}{210,210,210}
\definecolor{codekeyword}{RGB}{0,92,150}
\definecolor{codestring}{RGB}{160,60,0}
\definecolor{codecomment}{RGB}{95,95,95}
\definecolor{codeident}{RGB}{70,70,70}

\lstdefinestyle{reportstyle}{
    backgroundcolor=\color{codebg},
    frame=single,
    rulecolor=\color{codeframe},
    basicstyle=\ttfamily\small,
    keywordstyle=\color{codekeyword}\bfseries,
    stringstyle=\color{codestring},
    commentstyle=\color{codecomment}\itshape,
    identifierstyle=\color{codeident},
    showstringspaces=false,
    numbers=left,
    numberstyle=\tiny\color{gray},
    stepnumber=1,
    numbersep=8pt,
    tabsize=4,
    breaklines=true,
    breakatwhitespace=true,
    columns=fullflexible,
    captionpos=b,
    keepspaces=true,
    xleftmargin=1em,
    framexleftmargin=0.8em
}

\lstset{style=reportstyle}

\hypersetup{
    colorlinks=true,
    linkcolor=codekeyword,
    urlcolor=codekeyword,
    pdftitle={Manual de usuario del compilador de mini-lenguaje},
    pdfauthor={Ortega Novoa Octavio}
}

\begin{document}

\begin{titlepage}
\begin{center}

\begin{minipage}{0.49\textwidth}
\begin{flushleft}
\includegraphics[scale=0.12]{Imagenes/UNAM_logo.png}
\end{flushleft}
\end{minipage}
\begin{minipage}{0.49\textwidth}
\begin{flushright}
\includegraphics[scale=0.18]{Imagenes/LogoIngenieria.png}
\end{flushright}
\end{minipage}

\vfill
{\scshape\Large \textbf{UNIVERSIDAD NACIONAL AUTÓNOMA DE MÉXICO}} \\[4mm]
{\scshape\large \textbf{FACULTAD DE INGENIERÍA}} \\[8mm]

{\Large \textbf{Compiladores}} \\[4mm]
{\large Semestre 2026-2} \\[11mm]

{\Huge \textbf{Manual de usuario: Compilador de un mini-lenguaje con funciones, estructuras de control y expresiones booleanas}} \\[20mm]

\begin{tabular}{rl}
\textbf{Profesor:} & MANUEL ENRIQUE CASTAÑEDA CASTAÑEDA \\[2mm]
\textbf{Alumnos:} & Ortega Novoa Octavio \\
                  & [Nombre del segundo alumno] \\
                  & [Nombre del tercer alumno] \\[4mm]
\textbf{Fecha de entrega:} & 18 de marzo de 2026 \\
\end{tabular}

\end{center}
\end{titlepage}

\tableofcontents
\newpage

\section{Introducción}
Este manual explica cómo preparar el entorno, compilar el proyecto y ejecutar el compilador desde PowerShell en Windows. También resume la estructura de archivos, los requisitos de instalación y el formato general de uso del lenguaje de entrada.

El compilador fue desarrollado con Flex, Yacc/Bison y GCC. Su salida se organiza en archivos de texto dentro de \texttt{Codigo/output/}, lo que permite revisar tokens, variables, símbolos, código intermedio, código optimizado y código final de forma separada.

\section{Requisitos}
Antes de compilar el proyecto, el sistema debe contar con lo siguiente:
\begin{itemize}[label={-}]
    \item Windows con PowerShell.
    \item Flex instalado y disponible en la terminal.
    \item Yacc o Bison instalado y disponible como \texttt{yacc}.
    \item GCC instalado y disponible en la terminal.
    \item Permisos de escritura dentro de la carpeta del proyecto.
\end{itemize}

Si alguno de estos programas no está instalado, el compilador no podrá generarse correctamente.

\subsection{Verificación rápida}
Puedes comprobar que las herramientas estén disponibles con:
\begin{lstlisting}[language={}, caption={Verificación de herramientas}]
flex --version
yacc --version
gcc --version
\end{lstlisting}

Si alguno de esos comandos falla, primero instala o corrige la ruta de la herramienta correspondiente.

\section{Estructura del proyecto}
La carpeta principal del proyecto contiene estos elementos:
\begin{itemize}[label={-}]
    \item \texttt{Codigo/}: contiene el lexer, el parser, los archivos generados y los tests.
    \item \texttt{Codigo/lexer.l}: especificación léxica.
    \item \texttt{Codigo/parser.y}: gramática y acciones semánticas.
    \item \texttt{Codigo/tests/}: programas de prueba.
    \item \texttt{Codigo/output/}: archivos generados por el compilador.
    \item \texttt{Reporte	extunderscore Proyecto.tex}: reporte técnico del proyecto.
    \item \texttt{Manual	extunderscore Usuario.tex}: este manual.
\end{itemize}

\section{Compilación}
La compilación debe hacerse dentro de la carpeta \texttt{Codigo}. El flujo recomendado es exactamente el siguiente:

\begin{lstlisting}[language={}, caption={Compilación del compilador en PowerShell}]
PS C:\Users\Choi Beomgyu\Desktop\College\Proyecto LFyA - Compiladores\Codigo> flex lexer.l
PS C:\Users\Choi Beomgyu\Desktop\College\Proyecto LFyA - Compiladores\Codigo> yacc -y -d parser.y
PS C:\Users\Choi Beomgyu\Desktop\College\Proyecto LFyA - Compiladores\Codigo> gcc lex.yy.c y.tab.c -o compiler.exe
\end{lstlisting}

Si todo sale bien, se generará el ejecutable \texttt{compiler.exe} dentro de \texttt{Codigo}.

\subsection{Notas de compilación}
\begin{itemize}[label={-}]
    \item Debes ejecutar los comandos dentro de \texttt{Codigo}, no desde la raíz del proyecto.
    \item Si Flex o Yacc no están en el PATH, PowerShell mostrará un error de comando no reconocido.
    \item Si cambia el código fuente del lexer o del parser, es necesario volver a ejecutar los tres comandos anteriores.
\end{itemize}

\section{Ejecución}
Una vez creado el ejecutable, el compilador se corre con un archivo fuente del directorio \texttt{tests}. El ejemplo usado en este proyecto es:

\begin{lstlisting}[language={}, caption={Ejecución del compilador}]
PS C:\Users\Choi Beomgyu\Desktop\College\Proyecto LFyA - Compiladores\Codigo> .\compiler.exe tests/test5.c
\end{lstlisting}

Cuando la compilación es correcta, la terminal suele mostrar mensajes parecidos a estos:

\begin{lstlisting}[language={}, caption={Salida normal en terminal}]
[Info] Analisis lexico completado.
[Info] Analisis sintactico completado.
[Info] Tabla de tokens generada en output/tabla_tokens.txt
[Info] Tabla de variables generada en output/tabla_variables.txt
[Info] Tabla de simbolos generada en output/tabla_simbolos.txt
[Info] Codigo intermedio generado en output/codigo_intermedio.txt
[Info] Codigo final generado en output/codigo_final.txt
\end{lstlisting}

También puedes ejecutar los demás programas de prueba cambiando el nombre del archivo:
\begin{itemize}[label={-}]
    \item \texttt{tests/test1.c}
    \item \texttt{tests/test2.c}
    \item \texttt{tests/test3.c}
    \item \texttt{tests/test4.c}
    \item \texttt{tests/test5.c}
\end{itemize}

\subsection{Importante}
Los archivos de entrada deben escribirse con la sintaxis soportada por el compilador. Si el programa contiene errores léxicos o sintácticos, el compilador los reportará en consola y la salida generada puede quedar incompleta.

\section{Salida generada}
Cada ejecución produce archivos de texto en \texttt{Codigo/output/}. Los principales son:
\begin{itemize}[label={-}]
    \item \texttt{tabla\_tokens.txt}
    \item \texttt{tabla\_variables.txt}
    \item \texttt{tabla\_simbolos.txt}
    \item \texttt{codigo\_intermedio.txt}
    \item \texttt{codigo\_intermedio\_optimizado.txt}
    \item \texttt{codigo\_final.txt}
\end{itemize}

\subsection{Qué contiene cada archivo}
\begin{itemize}[label={-}]
    \item \texttt{tabla\_tokens.txt}: lista de tokens reconocidos por el lexer.
    \item \texttt{tabla\_variables.txt}: clasificación auxiliar de variables, cadenas y errores léxicos.
    \item \texttt{tabla\_simbolos.txt}: variables, funciones y estructuras detectadas semánticamente.
    \item \texttt{codigo\_intermedio.txt}: representación intermedia del programa.
    \item \texttt{codigo\_intermedio\_optimizado.txt}: versión simplificada del intermedio.
    \item \texttt{codigo\_final.txt}: traducción final a un pseudocódigo de ensamblador simple.
\end{itemize}

\subsection{Ejemplos de archivos}
Después de ejecutar \texttt{.\\compiler.exe tests/test5.c}, los archivos pueden verse de forma aproximada así:

\begin{lstlisting}[language={}, caption={Ejemplo de tabla de tokens}]
ID, main
LPAREN, (
RPAREN, )
LBRACE, {
KEYWORD, struct
ID, FlightControl
KEYWORD, switch
KEYWORD, case
KEYWORD, break
\end{lstlisting}

\begin{lstlisting}[language={}, caption={Ejemplo de tabla de variables}]
Nombre              Tipo      Categoria
missionPhase        char      variable
energyLevel         int       variable
message             string    cadena
tempValue           --        error
\end{lstlisting}

\begin{lstlisting}[language={}, caption={Ejemplo de tabla de símbolos}]
Nombre              Tipo      Categoria
manageFlight        void      funcion
computeEnergy       int       funcion
FlightControl       struct    estructura
PayloadBay          struct    estructura
\end{lstlisting}

\begin{lstlisting}[language={}, caption={Ejemplo de código intermedio}]
t1 = energyLevel + 5
ifFalse t1 goto L1
print "Estado nominal"
goto L2
L1:
    print "Revisar sistema"
L2:
\end{lstlisting}

\begin{lstlisting}[language={}, caption={Ejemplo de código final}]
MOV R1, energyLevel
ADD R1, 5
JZ L1
PRINT "Estado nominal"
JMP L2
L1:
PRINT "Revisar sistema"
L2:
\end{lstlisting}

\section{Cómo preparar un programa de entrada}
Los archivos de prueba están dentro de \texttt{Codigo/tests/}. El compilador acepta programas que usan:
\begin{itemize}[label={-}]
    \item Declaraciones de variables.
    \item Funciones.
    \item Estructuras \texttt{struct}.
    \item Condicionales \texttt{if} y \texttt{else}.
    \item Ciclos \texttt{while}, \texttt{do-while} y \texttt{for}.
    \item \texttt{switch/case/default/break}.
    \item \texttt{print} con expresiones y cadenas de texto.
\end{itemize}

\subsection{Ejemplo mínimo}
\begin{lstlisting}[language=C, caption={Programa de entrada mínimo}]
int main()
{
    int value = 10;
    print("Inicio del programa");
    print(value);
    return value;
}
\end{lstlisting}

\section{Consejos de uso}
\begin{itemize}[label={-}]
    \item Guarda el archivo fuente dentro de \texttt{Codigo/tests/} para evitar rutas largas.
    \item Recompila si cambias \texttt{lexer.l} o \texttt{parser.y}.
    \item Revisa \texttt{Codigo/output/} después de cada ejecución para validar los resultados.
    \item Si el compilador no arranca, confirma primero que \texttt{flex}, \texttt{yacc} y \texttt{gcc} respondan en la terminal.
\end{itemize}

\section{Resolución de problemas comunes}
\subsection{Flex o Yacc no se reconocen}
Si PowerShell muestra que \texttt{flex} o \texttt{yacc} no existen, significa que la herramienta no está instalada o no está agregada al PATH.

\subsection{GCC no genera el ejecutable}
Verifica que los archivos \texttt{lex.yy.c} y \texttt{y.tab.c} sí se generaron antes de ejecutar GCC. Si no existen, repite la compilación desde \texttt{flex} y \texttt{yacc}.

\subsection{No aparece la salida en output}
Revisa que el programa se haya ejecutado sin errores. Si el archivo de entrada tiene fallos, algunos archivos de salida pueden quedar vacíos o incompletos.

\section{Flujo recomendado}
El uso normal del proyecto es el siguiente:
\begin{enumerate}[label={\arabic*.}]
    \item Abrir PowerShell.
    \item Entrar a la carpeta \texttt{Codigo}.
    \item Ejecutar \texttt{flex lexer.l}.
    \item Ejecutar \texttt{yacc -y -d parser.y}.
    \item Ejecutar \texttt{gcc lex.yy.c y.tab.c -o compiler.exe}.
    \item Correr \texttt{.\\compiler.exe tests/test5.c} o el test que quieras revisar.
    \item Revisar los archivos generados en \texttt{Codigo/output/}.
\end{enumerate}

\section{Cierre}
Este manual resume los pasos necesarios para usar el compilador en Windows. Siguiendo el orden de compilación y ejecución indicado, cualquier usuario con Flex, Yacc y GCC instalados puede generar el ejecutable y probar los archivos de entrada del proyecto sin depender de instrucciones adicionales.

\end{document}
